% Compile with XeLaTeX. Template: moderncv classic (LPPL-1.3c).
\documentclass[11pt,a4paper,sans]{moderncv}
\moderncvstyle{classic}
\moderncvcolor{black}
\usepackage[UTF8,fontset=fandol]{ctex}
\usepackage[left=17mm,right=17mm,top=15mm,bottom=16mm]{geometry}
\setmainfont{texgyrepagella-regular.otf}[BoldFont=texgyrepagella-bold.otf,ItalicFont=texgyrepagella-italic.otf]
\setsansfont{texgyreheros-regular.otf}[BoldFont=texgyreheros-bold.otf,ItalicFont=texgyreheros-italic.otf]
\definecolor{color1}{HTML}{17675D}
\definecolor{color2}{HTML}{525754}
\colorlet{sectioncolor}{color1}
\colorlet{bodyrulecolor}{color1}
\setlength{\hintscolumnwidth}{25mm}
\setlength{\separatorcolumnwidth}{4mm}
\renewcommand*{\namefont}{\fontsize{29}{34}\mdseries\upshape}
\renewcommand*{\titlefont}{\fontsize{11}{15}\mdseries\upshape}
\renewcommand*{\sectionfont}{\fontsize{12}{16}\bfseries\upshape}
\name{欧阳}{Yang Ou}
\title{\texorpdfstring{高级算法工程师 · 哔哩哔哩\\内容安全 / VLM 后训练 / 多模态内容理解}{高级算法工程师 · 哔哩哔哩 · 内容安全 / VLM 后训练 / 多模态内容理解}}
\address{上海}{}{}
\email{ouyang980107@163.com}
\extrainfo{\href{https://orcid.org/0000-0002-1282-9022}{ORCID: 0000-0002-1282-9022}}
\begin{document}
\makecvtitle
\vspace{-4mm}
\small
\hypersetup{pdftitle={欧阳 Yang Ou | 内容安全与多模态算法},pdfauthor={欧阳 Yang Ou},pdfsubject={专业经历、代表项目、论文与专利}}
\fancyfoot[L]{\scriptsize\color{color2}Yang Ou / 内容安全与多模态研究}
\fancyfoot[R]{\scriptsize\color{color2}\thepage}

\section{专业定位}
\cvitem{}{专注\textbf{安全业务、VLM 后训练与细粒度内容理解}。在哔哩哔哩负责内容安全模型及审核多模态大模型能力，覆盖视频、图文、直播、社区与推荐场景。擅长将模糊业务规则转化为可学习的标签、训练目标和评估体系，从数据构建、模型训练到线上部署完整推进。持续将业务问题沉淀为研究课题，形成论文、专利与可复用的算法能力。}
\cvitem{科研积累}{\textbf{4 篇期刊论文、5 篇在投研究、5 项专利相关成果}；参与获得 \textbf{ICCV-CLVL 2025 Reliable VQA Challenge 第一名}。}

\section{核心技术能力}
\cvitem{后训练与推理}{监督微调（SFT）、强化学习与可验证奖励、On-Policy Distillation（OPD）、生成与分类联合训练、Adaptive CoT、推理成本优化。}
\cvitem{多模态理解}{图像、文本与上下文联合建模；细粒度多标签识别、风险严重度排序、结构化取证、幻觉控制与可靠视觉问答。}
\cvitem{数据与工程}{人机协同数据生产、困难样本挖掘、可控数据合成、模型蒸馏、召回与排序架构、训练及评估工具建设。}

\section{任职与代表项目}
\cventry{2022.06--至今}{高级算法工程师}{哔哩哔哩}{内容安全与多模态算法}{}{含 2022.06--2023.02 实习阶段；承担内容安全垂类模型与审核大模型能力建设，并协同推进跨业务落地及团队成员培养。}

\cvitem[.7em]{项目 01}{\textbf{严重度感知审核与自适应推理}\newline
\textbf{问题：}隐含风险依赖多元素共现与上下文，二分类难以表达连续的违规程度，高吞吐场景又限制大模型推理成本。\newline
\textbf{方法：}构建细粒度标签和严重度体系，以 SFT + RL 训练 VLM 排序能力；结合答案与过程奖励约束推理和结论的一致性。引入 Adaptive CoT，按样本难度分配推理开销；结合轻量召回与多模态排序架构上线。\newline
\textbf{产出：}在社区与推荐安全场景落地，并将严重度排序、过度思考问题分别沉淀为第一作者及共同第一作者研究。}

\cvitem[.7em]{项目 02}{\textbf{审核多模态大模型与结构化取证}\newline
\textbf{问题：}通用 VLM 缺少审核领域知识，风险信息覆盖与幻觉控制存在权衡，输出还需供不同下游策略稳定解析。\newline
\textbf{方法：}设计涵盖场景、主体、环境和细节的结构化取证模板；比较不同规模基础模型，采用 Caption with Box 监督微调与 OPD 蒸馏，增强风险感知、空间定位及输出稳定性。\newline
\textbf{产出：}从零搭建审核取证能力与评估迭代流程，支持下游按需使用证据，推进内容理解与处置决策的解耦。}

\clearpage
\cvitem{}{\textcolor{color1}{\Large 工程实践与技术沉淀}\newline\textcolor{color2}{从领域模型到数据、训练与线上系统}}
\section{更多代表项目}
\cvitem[.7em]{项目 03}{\textbf{色情低俗模型的多标签重构与跨业务复用}\newline
\textbf{方法：}将黑盒二分类规则拆解为通用多标签体系；联合优化 VLM 的生成与分类损失，以 Asymmetric Loss 和采样缓解长尾，并补充基于语言模型的文本风险识别能力。\newline
\textbf{产出：}统一视频、封面、图文及社区等场景的基础识别能力，支持可调节的召回及处置策略，提升维护效率和失败案例可归因性。}
\cvitem[.7em]{项目 04}{\textbf{图像检索与风险内容快速响应}\newline
\textbf{方法：}将图像检索升级为召回与精排两阶段架构；基于 Fast-ReID、对比学习与 ArcFace 度量学习优化局部图、模糊图识别。结合 PFC 加速训练，并利用 Louvain 社区发现与 DBSCAN 聚类构建千万级训练数据。引入样本入库标准、历史数据清洗与动态阈值，改善样本库质量和误召问题。\newline
\textbf{应用：}服务视频与图文风险内容检索，探索基于 Chinese CLIP 的文本检索图像能力，支持新风险类型的快速响应。}

\cvitem[.7em]{项目 05}{\textbf{直播场景识别与大模型蒸馏}\newline
\textbf{方法：}针对未成年人、血腥暴力等场景优化模型；重构年龄与未成年人识别方式，通过人脸检测、分类及卡通识别处理不同输入形态，减少非目标样本误召。将 VLM 的分类能力蒸馏到轻量视觉模型。\newline
\textbf{应用：}在直播推理吞吐约束下改善识别效果，兼顾审核效率与安全覆盖，并复用到相关场景识别任务。}

\cvitem[.7em]{项目 06}{\textbf{水印台标检测与可控样本合成}\newline
\textbf{方法：}比较并迭代 YOLO 检测方案，针对类别分散、标注不足以及裁剪和模糊输入构建可控数据生成引擎；通过样本分布配置与增强策略提升模型鲁棒性和扩类效率。\newline
\textbf{应用：}支持 90 余类水印台标识别，服务社区、视频、图文及版权相关业务。}

\cvitem[.7em]{项目 07}{\textbf{数据生产引擎与可复用训练工具}\newline
\textbf{数据侧：}搭建 Caption / QA 数据生产、人机协同细粒度标注与跨模型验证流程；围绕困难样本建立风格可控的生成和筛选流程，支持数据分布变化下的持续迭代。\newline
\textbf{训练侧：}开发 vision-classification 工具，集成模型骨干、损失函数、数据增强、自定义 Trainer、分类与蒸馏任务；沉淀离线评估和训练流程，减少重复开发与人工操作。}

\section{专利相关成果}
\cvitem{01 / 多模态}{一种结合多模态数据的审核优化方法。}
\cvitem{02 / 数据}{一种对用于训练审核模型的训练样本挖掘方法。}
\cvitem{03 / 部署}{一种针对多模态审核模型的蒸馏部署方法。}
\cvitem{04 / 审核}{一种基于违规程度的细粒度审核系统。}
\cvitem{05 / 可靠性}{一种基于自我反省和跨模型验证的多模态可靠问答方法。}

\section{协作与研究方法}
\cvitem{项目推进}{承担模型 Owner 职责，参与目标拆解、方案选择、实验设计、效果评估和上线协同；指导正式员工与实习生开展模型及数据工作。}
\cvitem{研究沉淀}{围绕真实业务约束提出问题，通过消融实验和困难样本分析验证方案，再沉淀为论文、专利、训练工具与技术文档。}

\clearpage
\cvitem{}{\textcolor{color1}{\Large 论文、竞赛与教育背景}\newline\textcolor{color2}{图学习基础 / 多模态后训练 / 可靠推理 / 内容安全}}
\section{已发表论文与公开预印本}
\cvitem[.6em]{2025 / 竞赛}{\textbf{Improving VQA Reliability: A Dual-Assessment Approach with Self-Reflection and Cross-Model Verification}\newline
Xixian Wu, \textbf{Yang Ou}, Pengchao Tian, Zian Yang, Jielei Zhang, Peiyi Li, Longwen Gao.\newline
\textbf{ICCV-CLVL 2025 Reliable VQA Challenge 第一名}；公开预印本 \href{https://arxiv.org/abs/2512.14770}{arXiv:2512.14770}。结合自我反省与跨模型验证提升回答可靠性。}
\cvitem[.6em]{2025 / 合作}{\textbf{Identification of perceptive users based on the graph convolutional network}\newline
Qiang Guo, Yong Luo, \textbf{Yang Ou}, Min Liu, Jian-Guo Liu.\newline
\textit{Expert Systems with Applications}. \href{https://doi.org/10.1016/j.eswa.2024.125844}{DOI: 10.1016/j.eswa.2024.125844}.}
\cvitem[.6em]{2022 / 一作}{\textbf{Identification of spreading influence nodes via multi-level structural attributes based on the graph convolutional network}\newline
\textbf{Yang Ou}, Qiang Guo, Jia-Liang Xing, Jian-Guo Liu.\newline
\textit{Expert Systems with Applications}. \href{https://doi.org/10.1016/j.eswa.2022.117515}{DOI: 10.1016/j.eswa.2022.117515}.}
\cvitem[.6em]{2022 / 一作}{\textbf{Identifying spreading influence nodes for social networks}\newline
\textbf{Yang Ou}, Qiang Guo, Jianguo Liu. \textit{Frontiers of Engineering Management}（在线发表）。\href{https://doi.org/10.1007/s42524-022-0190-8}{DOI: 10.1007/s42524-022-0190-8}.}
\cvitem[.6em]{2022 / 合作}{\textbf{Core-middle-periphery network model for China banking system}\newline
Na Chen, Jianguo Liu, Yihui Chen, Ding Tu, \textbf{Yang Ou}, Mingzhu Jiang.\newline
\textit{Frontiers in Physics}. \href{https://doi.org/10.3389/fphy.2022.981650}{DOI: 10.3389/fphy.2022.981650}.}

\section{在投研究}
\cvitem[.6em]{EMNLP 2026}{\textbf{Rethinking Content Moderation via Severity-Aware Ranking}\newline
\textbf{Yang Ou}, Jielei Zhang, et al. \textcolor{color1}{第一作者}。\href{https://bilibili-kuroko.github.io/}{KUROKO 项目主页}。}
\cvitem[.6em]{EMNLP 2026}{\textbf{When Not to Think: Mitigating the Overthinking Effect in Content Moderation}\newline
Yu Dai, \textbf{Yang Ou}, et al. \textcolor{color1}{共同第一作者}。}
\cvitem[.6em]{AAAI 2027}{\textbf{Zero-Shot Detection for AI-Generated Images via Uncertainty-Aware Manifold Deviation}\newline
Shutong Ding, \textbf{Yang Ou}, et al. \textcolor{color1}{共同第一作者}。}
\cvitem[.6em]{AAAI 2027}{\textbf{LatentLite: Lightweight Latent Reasoning with On-Demand Explanations}\newline
Bin Zhao, \textbf{Yang Ou}, et al. \textcolor{color1}{共同第一作者}。}
\cvitem[.6em]{ICLR 2027}{\textbf{Reconstruct, Then Evolve: Grounded Data Synthesis for Harmful Meme Detection}\newline
Bin Zhao, \textbf{Yang Ou}, Yihe Fan, Yuzhuo Zhang, Jielei Zhang, Junjie Wang, Qing Wang, Xudong Pan. \textcolor{color1}{第二作者}。}

\section{教育背景}
\cvitem{2020--2023}{\textbf{上海理工大学}，管理科学与工程硕士（2020.09--2023.03）。}
\cvitem{2016--2020}{\textbf{浙江万里学院}，电子商务学士（2016.09--2020.06）。}
\end{document}
